\chapter{Functional discovery, symbolic identification and estimating equations}
\label{v13:sw-functional}

\section{An operation is defined by its mathematical output}

A bank may ask a single software team to discover a graph, identify an effect
and estimate the value of an offer. Those requests have different outputs.
Discovery can return a selected directed graph, an equivalence class or a set
of candidates. Identification can return an expression in the observed law or
a nonidentification result. Estimation can return a number with a sampling
uncertainty claim. A shared product becomes misleading when it places all
three behind a field named ``causal effect'' without preserving these types.

Consider a renewal offer $A$, retained balance $Y$ and customer characteristics
$X$ available before the offer. A point-treatment question might be
\[
 \theta=\mathbb E\{Y(1)-Y(0)\}.
\]
Here the expectation is over a declared eligible population, not automatically
over the customers who accepted an offer. A discovery algorithm can help
propose relationships among $X,A,Y$, but a graph proposal does not by itself
establish that $Y(a)$ is identified. Even after identification, numerical
uncertainty depends on how customers were sampled and whether branch, household
or calendar dependence matters. Software comparison must preserve each stage.

The selected source paths considered here make these differences concrete.
DirectLiNGAM exposes a functional discovery construction; y0 exposes symbolic
identification; delicatessen exposes estimating equations; and statsmodels
exposes fitted treatment-effect operations with a joint GMM representation.
These are useful components of a workflow, with different scientific contracts.
The source inspections establish the specific interfaces and branches described
below. They do not establish execution speed, numerical agreement or empirical
coverage for this bank.\cite{v13_sw_lingam,v13_sw_y0,v13_sw_delicatessen,v13_sw_statsmodels}

\section{DirectLiNGAM: a selected order and the meaning of prior knowledge}

The inspected DirectLiNGAM class accepts a prior-knowledge matrix, an option to
apply it softly, a measure named \texttt{pwling}, \texttt{kernel} or
\texttt{pwling\_fast}, and a choice about adaptive-Lasso pruning. Its fit method
checks the prior matrix shape, selects variables sequentially, residualizes
remaining variables, stores a causal order and calls a separate adjacency
estimation routine. The fast branch refers to an additional GPU implementation;
the inspected dispatch is not a performance measurement.\cite{v13_sw_lingam}

The motivating functional family can be written
\[
 X_j=\sum_{k\in\operatorname{pa}(j)}b_{jk}X_k+\varepsilon_j,
\]
with an acyclic order and assumptions on the disturbances that distinguish this
family from unrestricted latent-confounded systems. A selected order narrows the
model much more than a CPDAG. That narrowing can be appropriate when its
functional assumptions are defensible; it cannot be justified by comparing only
the number of directed edges with a constraint-based method that deliberately
retains ambiguity. The present interface inspection does not verify the base
adjacency estimator or provide a finite confidence statement for the selected
order.

Prior knowledge is particularly consequential. The source describes zero as no
directed path, one as a directed path, and minus one as unspecified knowledge.
A bank statement that ``offer acceptance occurs after invitation'' is a temporal
constraint. It is not necessarily a statement that a directed path must exist.
Likewise, a field absent from the available dataset is not a justified absent
cause. An adapter should retain the original proposition, its variable order,
its path-versus-edge meaning and the evidence supporting it. Transposing an
array or treating an unknown entry as zero can change the scientific problem
without causing a matrix-shape error.

For funding analysis, useful adverse cases include a latent common market shock,
near-Gaussian disturbances and aggregation that mixes several decision times.
A comparison should ask whether the supplied functional model is valid in each
case, what output the package produces when it is not, and whether downstream
reporting preserves that limitation. It should not credit a plausible-looking
order as proof of an identified funding response. Constraint-based CPDAG/PAG
methods, alternative functional models and additional experimental information
are substantive alternatives, each requiring its own declared assumptions.

The appropriate first reuse candidate is an adapter around a reviewed operation.
A port would also require inspection of the inherited base class, total-effect
calculations, dependency terms and behavioral references. The observed MIT
license is useful evidence for that investigation; it is not a completed
dependency or maintenance assessment.\cite{v13_sw_lingam}

\section{y0: an expression can fail to exist before estimation begins}

The inspected \texttt{identify\_outcomes} function accepts an acyclic directed
mixed graph, treatment and outcome variables, optional conditions and an optional
ordering. With no conditions it dispatches to an ID routine; otherwise it
dispatches to IDC. It catches a dedicated \texttt{Unidentifiable} exception,
returning \texttt{None} by default and re-raising when \texttt{strict=True}.
This is a useful explicit boundary between a symbolic expression and a failure
to identify the requested query. It is only the dispatcher: the recursive
algorithms and graph validation need their own correctness review.\cite{v13_sw_y0}

For example, suppose the graph justifies adjustment for pre-offer $X$. The
corresponding functional, under the relevant causal assumptions, is
\[
 \theta=\int\{m_1(x)-m_0(x)\}\,dP_X(x),
 \qquad m_a(x)=\mathbb E(Y\mid A=a,X=x).
\]
The expression states what must be estimated. It supplies neither fitted
regressions nor a sampling distribution. If instead the graph contains an
unblocked latent common cause of $A$ and $Y$ and no identifying information,
another optimizer cannot manufacture this adjustment functional.

A product therefore needs separate outcomes for an identified expression, a
scientifically nonidentified query, an unsupported query language and an
execution error. Mapping all of them to a numeric missing value loses essential
information. Conversely, turning \texttt{None} into zero would state that the
effect is zero, which is a different scientific claim. The adapter should bind
the graph, query, conditioning variables and expression identities, and retain
the original failure type and diagnostic context.

Counterfactual coverage also needs care. The existence of an ID/IDC entry point
does not establish support for every nested counterfactual, path-specific effect
or transport query. The proper comparison has a query matrix: adjustment,
general intervention identification, conditional interventions and other
counterfactual families. Each cell should be tied to an inspected algorithm and
an independent small-graph reference. Package branding is not evidence for every
cell. A proposed symbolic adapter is justified as a candidate by this source;
general algorithmic correctness and comprehensive query support remain open.

\section{Stacked equations preserve the nuisance contribution}

Delicatessen's inspected \texttt{ee\_aipw} function returns an array of estimating
equations rather than a fitted effect object. Its parameter vector contains the
effect, two potential-outcome means and propensity/outcome coefficients. The
documentation and source identify a $(3+b+c)\times n$ array, where $b$ and $c$
are the numbers of propensity and outcome coefficients. The source selects
logistic or linear outcome regression from the outcome values unless the user
forces a continuous treatment of the outcome. It clips estimated propensity
probabilities only when bounds are supplied.\cite{v13_sw_delicatessen}

Writing $e(x)=P(A=1\mid X=x)$ and $m_a(x)=\mathbb E(Y\mid A=a,X=x)$, the
familiar augmented mean contribution has the form
\[
 \phi_a(O)=m_a(X)+\frac{\mathbf 1\{A=a\}}{p_a(X)}
                   \{Y-m_a(X)\},
 \qquad p_1=e,\quad p_0=1-e.
\]
This display describes the target contribution; its causal interpretation still
requires consistency, exchangeability and positivity for the declared target.
Stacking nuisance equations alongside equations for the two means makes their
joint role explicit. To obtain an estimator one must solve the system; to obtain
uncertainty one must review its Jacobian, covariance construction and sampling
unit. An array of equations is not an already qualified sandwich estimator.

This distinction matters when borrowers are clustered. If $\psi_i(\eta)$ is
an estimating-function contribution, an independent-customer covariance uses
individual outer products. A group construction instead uses sums
$\Psi_g(\eta)=\sum_{i\in g}\psi_i(\eta)$ and outer products of those group
sums, under an appropriate group sampling law. Neither replacing a variable
name nor using a generic robust standard error proves the required result for
shared calendars or persistent audit errors. The solver and covariance paths
must be inspected and qualified separately.

Truncation is also part of the method. Replacing $\widehat e$ by
$\max(\epsilon,\min(1-\epsilon,\widehat e))$ controls numerical weights but
changes the estimating equation. For a bank with almost no offers in a risk
stratum, clipping cannot create the missing counterfactual support. An adapter
must record the requested and applied bounds, how many observations were
affected, and the sensitivity of the answer. A stable answer after clipping is
not evidence that positivity held before clipping.

The same module exposes instrumental-variable and two-stage least-squares
estimating-equation entries. Their presence identifies a useful review route;
it does not establish weak-instrument inference, longitudinal g-estimation or
arbitrary dependent-data support. Reusing a selected equation as an independent
numerical reference may be valuable even when a complete estimator is not yet
an adoption candidate.\cite{v13_sw_delicatessen}

\section{statsmodels: population and covariance must survive the interface}

The inspected \texttt{TreatmentEffect} class describes five methods: IPW,
regression adjustment, AIPW, AIPW with weighted least squares and IPW with
regression adjustment. It states that inference uses a joint GMM representation
of the selection model, outcome model and effect functions. Its current outcome
model is documented as linear OLS. The result class exposes contrasts for an
effect and two potential-outcome means.\cite{v13_sw_statsmodels}

The population option distinguishes all, treated and untreated observations.
These correspond to different questions. With treatment propensity $e(X)$,
the all-population target averages over $P_X$, whereas an effect among treated
customers averages over $P_{X\mid A=1}$. Matching point estimates across packages
requires matching this measure as well as the outcome and treatment coding.
An apparent disagreement can be a population mismatch rather than a numerical
defect; apparent agreement in one synthetic sample can conceal the mismatch.

The class also exposes \texttt{ps\_bounds}, defaulting to $(0.001,0.999)$, and
validates strictly interior, ordered bounds. The selected source says the same
bounds apply to point estimates and GMM moment conditions. That is part of a
reproducible comparison, including a comparison with delicatessen's optional
truncation. Omitting these settings from a bridge specification means the two
programs need not be solving the same estimating equations.

One especially important detail is the internal \texttt{\_cov\_type} argument.
The inspected class states that it affects preliminary estimators but that final
GMM results correspond to HC0. A caller cannot infer clustered or calendar-robust
final inference merely from passing a covariance-related keyword. The early
module header also contains older prose about unavailable standard errors,
while the class and result implementation contain joint-GMM inference. Reviewing
the current operation resolves more than searching for one favorable or
unfavorable sentence.\cite{v13_sw_statsmodels}

For a banking baseline, this makes the candidate scope concrete: a declared
point-treatment population, linear outcome model, explicit propensity bounds
and a reviewed sampling law. Panel or time-series capabilities elsewhere in
statsmodels require separate inspection; they are not established by this one
module. The correct reuse decision is provisional numerical integration and
matched equation review, with refusal of unsupported dependence claims.

\section{A matched comparison must keep the question fixed}

A useful comparison first fixes the eligible population, intervention, observed
history and identification assumptions. It then asks each tool for its actual
output type: an order, symbolic expression, equation array or fitted estimate.
Small independent references can test the transitions. For discovery they
check order/edge conventions and unresolved alternatives. For identification
they check graph-query expression and failure semantics. For estimation they
check the target measure, nuisance fits, truncation and covariance construction.

Numerical disagreements should be classified by cause before they are ranked.
A different estimand is a semantic mismatch. A failed root solve is a numerical
failure. Nonidentification is a scientific outcome. An unsupported covariance
law is a scope limitation. Timing a successful branch while discarding the
others rewards a system for silently avoiding difficult requests. The full
comparison must retain them and distinguish source inspection from executed
behavior. These distinctions make the software survey useful for design and
investment even when it concludes that a component should be deferred.
